%% USPSC-Introducao.tex

% ----------------------------------------------------------
% Introdução (exemplo de capítulo sem numeração, mas presente no Sumário)
% ----------------------------------------------------------
\chapter[Introdução]{Introdução}
\label{Introdução}


O Algodão (\textit{Gossypium hirsutum L.}) é uma cultura com relevante importância no cenário internacional. 
Chegando a ser conhecido como "ouro branco" em diversos países, a fibra do algodão subsidia a indústria têxtil e tem um impacto
econômico de cerca de 600 bilhões de dólares ao ano \cite{Khan2020}. No Brasil, a produção de pluma de
algodão gira em torno de 2,5 milhões de toneladas ao ano, em uma área plantada com cerca de 1,6 milhão
de hectares \cite{conab2024}.

Uma importante etapa dentro cadeia produtiva do algodão é a classificação da qualidade da fibra de
acordo com níveis internacionais. Esses níveis de qualidade determinam o preço de comercialização e o
poder de negociação do produto \cite{yasar2021quality}.

A classificação internacional da fibra começou com os americanos, em 1906, pelo Departamento de
Agricultura dos Estados Unidos (USDA), que estabeleceu a padronização da qualidade da fibra, em um
sistema que existe até hoje. Atualmente o USDA estabelece 15 padrões de classificação de algodão, que
são atualizados todos os anos e servem como referência para as classificações visuais realizadas \cite{cottoninc2017}.

No Brasil, o processo de classificação é regulamentado pela Instrução Normativa n. 24/2016 do
Ministério de Estado da Agricultura, Pecuária e Abastecimento (MAPA), que traça padrões para a
classificação, com requisitos de identidade e qualidade, amostragem e apresentação do produto \cite{mapa2016}.

No processo produtivo do algodão a classificação é feita por amostragem, e pode ser realizada de duas
formas: por método visual ou por método instrumental. A classificação visual é realizada por classificador
(pessoa física) tendo como base Padrões Físicos Universais, levando em conta a cor das fibras, a presença
de folhas que irá caracterizar as impurezas, as contaminações de matérias estranhas e o modo de
preparação (beneficiamento) do produto. Já a classificação instrumental é realizada por equipamento do
tipo HVI (\textit{High Volume Instrument}) \cite{mapa2016}. A classificação instrumental demanda um tempo de análise 
considerável, bem como custos elevados, preparo das amostras e condições controladas em laboratório \cite{Oliveira2022}.

Após a produção do fardo do algodão beneficiado e a classificação é realizada a etapa de "emblocamento"
que consiste na criação de blocos de algodão de aproximadamente 25 mil toneladas, compostos por fardos
que possuem a mesma classificação (geralmente entre 110 e 125 fardos). Cada bloco é a menor unidade
de comercialização.

Como o processo de beneficiamento nas algodoeiras ocorre 24 horas por dia e 7 dias por semana e devido
a limitações logísticas, existe a necessidade de se realizar o "emblocamento" no menor tempo possível.
Por isso, as empresas optam por realizar esta etapa do processo de acordo com a classificação visual, já
que os resultados da classificação por instrumentos podem demorar até 24 horas para chegar.

Mesmo com a classificação visual, que é mais rápida que a por instrumentos, existe um ponto de lentidão
no processo, pois é necessário que as amostras sejam enviadas para as unidades de classificação e
somente depois que o emblocamento é realizado. Ter as unidades de classificação dentro da indústria não
se mostra viável pela necessidade de se ter um ambiente controlado para garantia da qualidade das
amostras. Assim, ter equipamentos que possam ser instalados na indústria (após a geração dos fardos),
que possam coletar, processar e analisar de forma automatizada as amostras de cada fardo, se mostra uma
alternativa com grande potencial para melhoria do processo. Neste cenário, um sistema de visão
computacional parece ser o mais adequado.

Visão computacional é a ciência responsável pela visão de uma máquina, pela forma como um
computador enxerga o meio à sua volta, extraindo informações significativas a partir de imagens
capturadas por diversos dispositivos como câmeras, sensores, scanners, etc. Estas informações permitem
reconhecer, manipular e pensar sobre os objetos que compõem uma imagem \cite{ballard1982computer}.

Por volta da década de 70, foram iniciados os primeiros estudos na área de visão computacional. Ao longo
desses anos uma série de pesquisas foram conduzidas e técnicas de análise foram elaboradas,
principalmente com o uso de inteligência artificial. Assim, a visão computacional fornece ao computador
uma infinidade de informações precisas a partir de imagens e vídeos, de forma que o computador consiga
executar tarefas inteligentes, simulando e aproximando-se da inteligência humana \cite{oliveira2024visao}.

Um sistema básico de visão computacional consiste em quatro características: iluminação, câmera,
computador e software de processamento de imagens. Esse sistema vai permitir a captura da imagem, sua
conversão digital para armazenamento e processamento no computador e o software de processamento de
imagens será responsável pelo aprimoramento, segmentação, detecção ou classificação da imagem
\cite{zhang2014computer}

O sistema de iluminação deve ser capaz de destacar as características de interesse do objeto, em
detrimento das demais características. É um sistema de extrema importância para se garantir a eficiência e
a acurácia de um sistema de visão computacional, além de ter uma grande influência na qualidade das
imagens. As chaves de um bom sistema de iluminação são aquelas que provêm uma iluminação uniforme
com um espectro de características e que produzem imagens de alta qualidade \cite{thomasson2005}.

A câmera é o componente principal do sistema de visão computacional, que é usado para a aquisição da
imagem e para a troca de dados com o computador. Existem diferentes tipos de dispositivos que capturam
imagens em diversos espectros como monocromáticas, RGB, ultrassom, raios-x, NIR, etc \cite{brosnan2002}.

A classificação visual por se tratar de um processo realizado manualmente e, apesar de possuir um
padrão, o resultado está sujeito à interpretação a à subjetividade do profissional que a está executando,
muitas vezes abrindo margens para questionamentos e rejeições na fase de negociação da mercadoria.

Diante deste cenário, justifica-se a construção de um sistema que possibilite realizar a classificação visual
do produto, seguindo todas as regulamentações necessárias e que elimine a subjetividade do processo e
com velocidade. Assim, propõe-se este trabalho com o objetivo de construir um sistema automatizado de
classificação do algodão, utilizando modelos de visão computacional e inteligência artificial, com o
intuito de se ter uma maior padronização das classificações realizadas, bem como reduzir o número de
rejeições da fase de negociação da mercadoria. Objetivos específicos: O hardware deve possibilitar a
coleta e a classificação da amostra, sem depender de outros equipamentos ou softwares; A velocidade
para a coleta da imagem, processamento e inferência deve ser adequada à demanda que a operação exige.
Deve ser possível integrar via API os dados das amostras e suas classificações.